\documentclass{ctexart}
\usepackage{microtype}
\usepackage[
    a5paper,
    inner=2.2cm,
    outer=1.5cm,
    top=1.8cm,
    bottom=1.8cm]{geometry}
\setlength{\parskip}{0.5em}
\setlength{\parindent}{0pt}
\usepackage{setspace}
\setstretch{1.7}
\usepackage{xpinyin}
\xpinyinsetup{ratio={.5}, format={\scriptsize}, vsep={-0.5em}, hsep={0.3em}}
\AtBeginDocument{% Custom pinyin for Heart Sutra (心經)
% Buddhist terms with special pronunciations

% 般若 - Prajna (bō rě, not pán ruò)
\setpinyin{般}{bō}
\setpinyin{若}{rě}

% 波羅蜜多 - Paramita
\setpinyin{羅}{luó}

% 舍利子 - Sariputra
\setpinyin{舍}{shè}

% 菩提薩埵 - Bodhisattva
\setpinyin{薩}{sà}
\setpinyin{埵}{duǒ}

% 涅槃 - Nirvana
\setpinyin{槃}{pán}

% 阿耨多羅三藐三菩提 - Anuttara-samyak-sambodhi
\setpinyin{耨}{nòu}
\setpinyin{藐}{miǎo}

% 罣礙 - hindrance
\setpinyin{罣}{guà}
\setpinyin{礙}{ài}

% Mantra: 揭諦揭諦 波羅揭諦 波羅僧揭諦 菩提薩婆訶
\setpinyin{揭}{jiē}
\setpinyin{諦}{dì}
\setpinyin{婆}{pó}
\setpinyin{訶}{hē}
}

\setCJKmainfont{AR PL KaitiM Big5}
\title{\Huge 心經}
\author{}
\date{}

\pagestyle{plain}

\begin{document}

\maketitle
\vspace{1em}
\begin{center}
般若波羅蜜多心經\\[0.5em]
Prajñāpāramitāhṛdaya
\end{center}
\vfill
\begin{center}
cbeta: T0251
\end{center}
\thispagestyle{empty}
\clearpage

\Large
\xpinyin*{觀自在菩薩，行深般若波羅蜜多時，照見五蘊皆空，度一切苦厄。}

\xpinyin*{舍利子，色不異空，空不異色；色即是空，空即是色。}

\xpinyin*{受想行識，亦復如是。}

\xpinyin*{舍利子，是諸法空相，不生不滅，不垢不淨，不增不減。}

\xpinyin*{是故空中無色，無受想行識。}

\xpinyin*{無眼耳鼻舌身意，無色聲香味觸法。}

\xpinyin*{無眼界，乃至無意識界。}

\xpinyin*{無無明，亦無無明盡，乃至無老死，亦無老死盡。}

\xpinyin*{無苦集滅道，無智亦無得。}

\xpinyin*{以無所得故，菩提薩埵，依般若波羅蜜多故，心無罣礙。}

\xpinyin*{無罣礙故，無有恐怖，遠離顛倒夢想，究竟涅槃。}

\xpinyin*{三世諸佛，依般若波羅蜜多故，得阿耨多羅三藐三菩提。}

\xpinyin*{故知般若波羅蜜多，是大神咒，是大明咒，是無上咒，是無等等咒，能除一切苦，真實不虛。}

\xpinyin*{故說般若波羅蜜多咒，即說咒曰：}

\vspace{1em}
\begin{center}
\Large
\xpinyin*{揭諦揭諦，}\\
\xpinyin*{波羅揭諦，}\\
\xpinyin*{波羅僧揭諦，}\\
\xpinyin*{菩提薩婆訶。}
\end{center}
\vspace{1em}

\clearpage
\null
\thispagestyle{empty}

\end{document}
